\documentclass{tietreport}

% Import images
\usepackage{graphicx}
\graphicspath{{../}} % existing project diagrams live in docs/

% Bibliography configuration
\usepackage[
  date=year,
  style=alphabetic,
  backend=biber,
  sorting=nty,
  maxnames=5,
  minnames=3,
  maxalphanames=4,
  minalphanames=3,
  backref=true,
  doi=false,
  isbn=false,
  url=false,
  eprint=false
]{biblatex}
\addbibresource{references.bib}

%% ==================================================
%% CUSTOMIZE THESE FIELDS FOR YOUR PROJECT
%% ==================================================

% Course/Lab header
\TitlePageHeader{%
  UCS503P - Software Engineering Project%
}

% Main project title
\ProjectTitle{%
  CVATS -- AI-Powered ATS Resume Scanner%
}

% Subtitle or project description
\ProjectSubTitle{%
  Job-Description-Based Resume Matching, Gap Analysis,
  Resume Enhancement and Mock Interview System%
}

% Student information (up to 4 students)
\author{%
  \texttt{1024030509} Yashasvi Ranjan \\
  \texttt{1024030508} Kamesh Yadav \\
  \texttt{1024030499} Taniksha \\
  \texttt{1024030521} Anshika%
}

% Group number, degree, year, program
\TitlePageSubText{%
  Group: \texttt{CVATS} BE Third Year, CSE%
}

% Faculty advisor/guide name
\AdvisorName{%
  Dr. Nisha Thakur%
}

% Evaluation period and department info
\TitlePageFooterText{{%
    \large Mid-Semester Evaluation \\[0.2cm]
    \bfseries Computer Science and Engineering
    Department%
  } \\[0.15cm]
  Thapar Institute of Engineering and Technology,
  Patiala%
}

%% ==================================================

\begin{document}

% Generate title page
\maketitle

% Table of contents (auto-generated from chapters)
\tableofcontents

% ===== CHAPTER 1: INTRODUCTION =====

\chapter{Introduction}

\section{Background and Context}

Applicant tracking systems (ATS) are used by most employers to filter
resumes before a human recruiter reads them. Job seekers rarely see how
their resume is judged, and a well-qualified candidate can be rejected
because skills are phrased differently from the job description (JD) or
because relevant experience is buried. Existing resume checkers are
usually opaque: they return a single number with no evidence, and
generative tools may ``improve'' a resume by inventing skills or
achievements the candidate never had.

CVATS is a tool for \emph{job seekers}, not a recruiter or HR screening
platform. A user uploads a resume, pastes a job description, and the
system walks them through resume parsing, JD parsing, an explainable
match score, gap analysis, grounded resume enhancement, a role-specific
mock interview, and a stored history of all results~\cite{cvatsrepo}. The
match score is an estimate and is not an official score from any
employer's ATS.

\subsection{Problem Statement}

Job seekers lack a transparent, trustworthy way to understand how well a
resume matches a specific job description, what exactly is missing, and
how to improve and rehearse. Current tools either give unexplained scores
or use large language models (LLMs) in ways that can fabricate content.
The problem tackled here is to build a deployable web system that
(i) computes a reproducible, evidence-backed match score,
(ii) separates \emph{weakly presented} content from \emph{genuinely
missing} content, (iii) rewrites resume text without inventing facts, and
(iv) prepares the user for interviews, while degrading honestly when the
AI provider is unavailable.

\section{Project Objectives}

\begin{enumerate}
  \item \textbf{Structured parsing:} Extract text from PDF/DOCX resumes and
    pasted JDs into structured, user-correctable data in which every
    parsed item traces back to the source text.
  \item \textbf{Explainable matching:} Produce a deterministic, weighted
    match score over skills, experience, keywords, education and projects,
    with evidence, deductions and confidence per category.
  \item \textbf{Gap analysis:} Classify findings as \emph{Things to Improve}
    or \emph{Things to Learn/Add}, with priorities, suggested actions and
    saved progress.
  \item \textbf{Grounded enhancement:} Generate reworded resume content
    validated against the original so that invented skills, numbers,
    employers, titles, dates and credentials are reverted.
  \item \textbf{Mock interview:} Generate role-specific questions from the
    resume, JD and gaps, evaluate answers as coaching feedback, and store a
    final report.
  \item \textbf{Secure and deployable:} Provide authenticated, owner-scoped
    access with a REST API, documented with Swagger, and a deployable
    React frontend and Node.js backend.
\end{enumerate}

These objectives were defined to be specific and measurable: each maps to
a concrete API area and database entity that can be demonstrated to
evaluators.

% ===== CHAPTER 2: METHODOLOGY & DESIGN =====

\chapter{Methodology and Design}

\section{System Architecture}

CVATS follows a three-tier architecture: a single-page React client, a
stateless Express REST API, and a PostgreSQL database accessed through the
Prisma ORM. All LLM calls are made only from the backend, behind a
provider abstraction, so API keys never reach the browser.
Figure~\ref{fig:dfd0} shows the context-level data flow and
Figure~\ref{fig:component} the component structure.

\begin{figure}[h]
\centering
\includegraphics[width=0.85\textwidth]{"DFD_0-Level DFD.drawio.png"}
\caption{Level-0 data flow diagram of CVATS}
\label{fig:dfd0}
\end{figure}

\begin{figure}[h]
\centering
\includegraphics[width=0.75\textwidth]{"component_diagram.drawio.png"}
\caption{Component diagram}
\label{fig:component}
\end{figure}

\subsection{High-Level Design}

The user registers and logs in (JWT), uploads a resume and submits a JD.
Both are parsed into structured JSON validated by Zod schemas. Starting an
analysis creates a record with status \texttt{PENDING}; a background
runner moves it through \texttt{PROCESSING} to \texttt{COMPLETED} or
\texttt{FAILED}, and the client polls for progress. The deterministic
scoring engine computes the match; recommendations, enhancement versions
and interview sessions are generated from that analysis and persisted.
Figures~\ref{fig:class} and~\ref{fig:sequence} give the class structure and
the main interaction sequence, and Figure~\ref{fig:deploy} the deployment
view.

\begin{figure}[h]
\centering
\includegraphics[width=0.7\textwidth]{"class_diagram.drawio.png"}
\caption{Class diagram}
\label{fig:class}
\end{figure}

\begin{figure}[h]
\centering
\includegraphics[height=0.6\textheight]{"Sequence Diagram.drawio.png"}
\caption{Sequence diagram of the analysis workflow}
\label{fig:sequence}
\end{figure}

\begin{figure}[h]
\centering
\includegraphics[width=\textwidth]{"cvats_deploymentDiagarm.png"}
\caption{Deployment diagram}
\label{fig:deploy}
\end{figure}

\subsection{Technical Stack}

\begin{itemize}
  \item \textbf{Framework/Language:} Frontend in React 18 with Vite,
    Tailwind CSS and React Router; backend in Node.js, Express and
    TypeScript, with Zod for validation.
  \item \textbf{Database:} PostgreSQL with Prisma ORM and additive
    migrations; entities include User, Resume, JobDescription,
    JobAnalysis, AnalysisRecommendation, ResumeEnhancement and the
    interview session, question, answer and evaluation tables.
  \item \textbf{Tools:} Git, Swagger UI for API documentation, Jest for
    tests, and Vercel configuration for frontend deployment.
  \item \textbf{Libraries:} \texttt{pdf-parse} (PDF), \texttt{mammoth}
    (DOCX), \texttt{bcryptjs}, \texttt{jsonwebtoken}, \texttt{helmet},
    \texttt{express-rate-limit}, \texttt{multer}. The LLM is Groq
    (default model \texttt{openai/gpt-oss-120b}) behind an abstraction that
    also supports OpenAI and OpenAI-compatible servers such as
    Ollama~\cite{groqdocs,prisma}.
\end{itemize}

\section{Additional Design Diagrams}

The following diagrams from the design phase complete the model of the system.

\begin{figure}[h]
\centering
\includegraphics[width=0.85\textwidth]{"Use case Diagram.drawio.png"}
\caption{Use case diagram}
\label{fig:usecase}
\end{figure}

\begin{figure}[h]
\centering
\includegraphics[height=0.6\textheight]{"ACTIVITY DIAGRAM .drawio.png"}
\caption{Activity diagram}
\label{fig:activity}
\end{figure}

\begin{figure}[h]
\centering
\includegraphics[width=0.85\textwidth]{"SWIMLANE.drawio.png"}
\caption{Swimlane diagram}
\label{fig:swimlane}
\end{figure}

\begin{figure}[h]
\centering
\includegraphics[width=0.85\textwidth]{"Collaboration Diagram.drawio.png"}
\caption{Collaboration diagram}
\label{fig:collab}
\end{figure}

\begin{figure}[h]
\centering
\includegraphics[width=0.85\textwidth]{"DFD_1-Level DFD.drawio.png"}
\caption{Level-1 data flow diagram}
\label{fig:dfd1}
\end{figure}

\begin{figure}[h]
\centering
\includegraphics[width=0.85\textwidth]{"DFD_2-Level DFD.drawio.png"}
\caption{Level-2 data flow diagram}
\label{fig:dfd2}
\end{figure}

\section{Implementation Details}

\subsection{Module 1: Parsing and Match Scoring}

Resume text is extracted from PDF or DOCX (image-only PDFs are not
supported, as there is no OCR) and converted to a structured resume
(skills, experience, projects, education, certifications). The JD parser
extracts required versus preferred skills, experience and education, with
no hardcoded default requirements; the user may correct either result.

The scoring engine is deterministic and reproducible. Each requirement is
looked up in a skill taxonomy with normalisation and related-skill
matching; the strongest single piece of evidence counts once, so
repetition cannot inflate the score. Credit depends on evidence strength
(demonstrated $=1$, claimed $=0.85$, listed $=0.7$), required items weigh
three times preferred ones, and related skills earn half credit. The LLM
may only help classify uncertain matches, and only when it cites a real
resume snippet. The default category weights are in
Table~\ref{tab:weights} and are normalised automatically.

\begin{table}[h]
\centering
\begin{tabular}{|l|r|}
  \hline
  \textbf{Category} & \textbf{Default Weight} \\
  \hline
  Skills & 0.35 \\
  \hline
  Experience & 0.25 \\
  \hline
  Keywords & 0.15 \\
  \hline
  Projects & 0.15 \\
  \hline
  Education & 0.10 \\
  \hline
\end{tabular}
\caption{Default match-score weights}
\label{tab:weights}
\end{table}

\subsection{Module 2: Gap Analysis, Enhancement and Mock Interview}

\textbf{Gap analysis} derives recommendations from the match result and
separates \emph{Things to Improve} (present but weakly presented) from
\emph{Things to Learn/Add} (genuinely missing), each with a priority,
evidence, suggested actions and learning topics; progress is saved.

\textbf{Resume enhancement} asks the LLM to reword content using only the
original resume. A grounding validator compares the output with an index
of the source and reverts invented skills, numbers, employers, titles,
dates and credentials. Users get a side-by-side comparison, editing,
versioning, approval and TXT/DOCX download; the original upload is never
modified.

\textbf{Mock interview} generates questions from the resume, JD and gaps,
evaluates each answer as coaching feedback (not a hiring decision), allows
sessions to be resumed, and stores a final report.

\subsection{Module 3: Security and Reliability}

Resumes, JDs and interview answers are treated as untrusted data and
wrapped so embedded instructions are not followed. Scores and evaluations
are written only by trusted server code. LLM output is validated against
schemas, malformed output gets at most one repair attempt, retries are
bounded, and calls are limited by a concurrency semaphore and queue.
Passwords are hashed with bcrypt, routes are protected by JWT and scoped to
the resource owner, and rate limiting applies per IP and per user on
expensive AI endpoints. When the provider is unavailable the analysis is
recorded as \texttt{FAILED}; fabricated results are never returned.

% ===== CHAPTER 3: RESULTS & EVALUATION =====

\chapter{Results and Evaluation}

\section{Testing Methodology}

\begin{itemize}
  \item \textbf{Unit Testing:} Jest with ts-jest for the deterministic
    components such as the scoring engine, weight normalisation, skill
    normaliser and grounding validator.
  \item \textbf{Integration Testing:} Jest integration tests against a
    separate test database (\texttt{.env.test}) with Prisma migrations
    applied, exercising the REST endpoints.
  \item \textbf{Performance Testing:} Bounded LLM concurrency, queueing
    and timeouts are configurable; formal load testing has not yet been
    carried out.
  \item \textbf{User Testing:} Manual end-to-end walkthroughs of the
    complete flow, using the demo mode (\texttt{VITE\_DEMO\_MODE}) with
    sample data, isolated from real workflows.
\end{itemize}

\section{Performance Metrics}

No quantitative benchmark figures have been recorded yet, so none are
reported. Table~\ref{tab:status} instead summarises the delivery status of
each objective against the implemented system.

\begin{table}[h]
\centering
\begin{tabular}{|l|l|l|}
  \hline
  \textbf{Objective} & \textbf{Evidence in System} &
  \textbf{Status} \\
  \hline
  Structured parsing & Resume and JD parse endpoints & Implemented \\
  \hline
  Explainable matching & Deterministic engine, per-category evidence
  & Implemented \\
  \hline
  Gap analysis & Improve vs.\ Learn/Add recommendations & Implemented \\
  \hline
  Grounded enhancement & Validator, versioning, DOCX/TXT export
  & Implemented \\
  \hline
  Mock interview & Questions, evaluation, report & Implemented \\
  \hline
  Secure deployment & JWT, rate limits, Swagger, Vercel config
  & Implemented \\
  \hline
\end{tabular}
\caption{Objective delivery summary}
\label{tab:status}
\end{table}

\section{Analysis}

\begin{itemize}
  \item \textbf{Objectives:} All six objectives have a working
    implementation in the backend and frontend.
  \item \textbf{Strengths:} Reproducible scoring and the grounding
    validator address the two main weaknesses of existing tools, opacity
    and fabrication.
  \item \textbf{Challenges:} LLM output can be malformed or unavailable,
    and resumes are untrusted text that may contain injected
    instructions.
  \item \textbf{Resolution:} Schema validation with a single repair
    attempt, bounded retries, explicit \texttt{FAILED} states, and
    wrapping of untrusted input.
\end{itemize}

% ===== CHAPTER 4: CONCLUSION =====

\chapter{Conclusion}

\section{Summary of Work}

We designed and built CVATS, a full-stack system that guides a job seeker
from resume and JD parsing through explainable scoring, gap analysis,
grounded enhancement and mock interviews, with history and restorable
results. The implementation meets the objectives set in Chapter~1 at the
functional level.

\section{Key Findings}

\begin{itemize}
  \item \textbf{Finding 1:} Keeping scoring deterministic and using the LLM
    only as a verified assistant gives reproducible, explainable scores.
  \item \textbf{Finding 2:} A post-generation grounding validator is
    necessary to stop LLM rewrites from inventing credentials or figures.
  \item \textbf{Finding 3:} Honest failure handling matters: reporting a
    provider outage is preferable to returning fabricated output.
\end{itemize}

\section{Future Work and Recommendations}

\begin{enumerate}
  \item OCR support for image-only PDF resumes~\cite{prisma}
  \item Calibrating scoring weights against labelled resume--JD pairs
  \item Background job queue and caching for scalability
  \item Browser extension or job-board integration to import JDs
  \item Formal load testing and production monitoring
\end{enumerate}

\section{Final Remarks}

Building CVATS showed that trustworthy AI assistance depends as much on
validation, evidence and failure handling as on the model itself. The
system gives job seekers a transparent way to understand and improve their
fit for a role.

% ===== BIBLIOGRAPHY =====

% Print bibliography with heading in TOC
\printbibliography[heading=bibintoc]

\end{document}
